% Generated draft: requires mathematical and source review.
\documentclass[11pt]{article}
\usepackage[margin=25mm]{geometry}
\usepackage{fontspec}
\setmainfont{Latin Modern Roman}
\usepackage{amsmath,amssymb,mathtools,unicode-math}
\setmathfont{Latin Modern Math}
\usepackage{enumitem,needspace,xurl,hyperref}
\setlength{\parindent}{0pt}
\setlength{\parskip}{5pt}
\setlength{\emergencystretch}{3em}
\newcommand{\N}{\mathbb{N}}
\newcommand{\Z}{\mathbb{Z}}
\newcommand{\Q}{\mathbb{Q}}
\newcommand{\R}{\mathbb{R}}
\newcommand{\C}{\mathbb{C}}
\newcommand{\F}{\mathbb{F}}
\newcommand{\A}{\mathbb{A}}
\newcommand{\PP}{\mathbb{P}}
\newcommand{\E}{\mathbb{E}}
\newcommand{\eps}{\varepsilon}
\newcommand{\dd}{\,\mathrm{d}}
\begin{document}
\section*{UnsolvedMath ID: 30000882}
\textbf{Source problem code:} catalogue entry
\subsection{Definitions and mathematical statement}
The source bundles four separate group-ring questions.
\begin{enumerate}
\item For finite groups $G,H$, if $Z(\Z[G])\cong Z(\Z[H])$ as rings, must the groups have equivalent character tables, with class sums corresponding? A class sum is $\sum_{g\in C}g$ for a conjugacy class $C$; these give the usual integral basis of the center. Ordinary character tables have irreducible complex characters as rows and conjugacy classes as columns. The source does not define table equivalence or say whether the given ring isomorphism must preserve class sums, or whether some suitable correspondence is sought. These are distinct requirements.
\item For torsion-free $G$, is $U(K[G])$ trivial? The field $K$ is unspecified. Literally a trivial unit group would mean only the identity, whereas group elements and nonzero scalars already give units. The usual intended ``trivial units'' property is $U(K[G])=\{ag:a\in K^\times,g\in G\}$. This likely convention is recorded for review, not silently substituted for the ambiguous wording.
\item If the normalizer problem holds for every finite normal subgroup of $G$, must it hold for $G$, and can groups satisfying it be classified? In the common integral-group-ring convention the property is
\[
 N_{U(\Z[G])}(G)=G\,Z(U(\Z[G])).
\]
The source does not specify this ring convention, or whether the premise concerns each finite normal subgroup's own group ring or its relative normalizer inside the ambient unit group. No one interpretation is selected as a definitive universal theorem.
\item For the finite-dimensional algebra $A$ arising from a noncentral finite-conjugacy center of the unit group of a $\Z$-order, determine its structure and Wedderburn--Malcev decomposition. The finite-conjugacy center of a group $U$ is $\{u\in U:\text{the conjugacy class of }u\text{ in }U\text{ is finite}\}$. The record does not define the construction of $A$ from that subgroup or the order. A Wedderburn--Malcev description ordinarily separates the Jacobson radical from a suitable semisimple subalgebra, but the field, hypotheses, and desired classification data are absent.
\end{enumerate}

All four questions are retained without imposing one common finite-group hypothesis. The fourth concerns orders and a derived algebra rather than automatically the group algebra from the first three. The dated assessment counts three themes, whereas the statement and verification explicitly contain four; it cannot supply a status assessment of the fourth by omission.
\subsection{Short English statement}
The entry contains four distinct questions about center isomorphisms, trivial units, normalizer properties, and an algebra associated with finite conjugacy classes in order-unit groups. Several terms and scopes are ambiguous, especially trivial units and the normalizer premise.
\subsection{Sources}
\begin{itemize}
\item Local UnsolvedMath registry, permanent ID 30000882, code catalogue entry, ``Centers, Normalizers, and Units of Integral Group Rings.'' Proposed year 2007; proposer unspecified.
\item \emph{Mini-Workshop: Arithmetik von Gruppenringen} (2008), Oberwolfach Reports, pp.\ 3209--3240, DOI 10.4171/owr/2007/55; \url{https://doi.org/10.4171/owr/2007/55}. The displayed report year 2008 differs from the DOI and proposed year 2007.
\item Supplied ``normalizer problem results'': \url{https://doi.org/10.1006/jabr.2001.8724}.
\item Supplied ``centers of integral group rings'': \url{https://arxiv.org/abs/math/0612436}.
\item Statement status corrected\_verified; the supplied note explicitly preserves all four questions and restores integral-ring/algebra notation.
\item Literature triage checked 2026-08-21: partially solved. It reports substantial special cases for character-table correspondence, torsion-free trivial-unit questions, and the normalizer property, but no general resolution of those three themes; it does not separately assess the fourth algebra question. Created 2026-08-25; updated 2026-09-16. Exact hypotheses and current statuses require independent verification.
\item Registry curation by Ulam AI / UnsolvedMath Contributors, CC BY 4.0; underlying report and literature provenance remain as supplied. Unreviewed statement-only draft.
\end{itemize}
\end{document}
